% ============================================================================
% 05_ers_ieee830.tex — Cuerpo del entregable "ERS IEEE Std 830-1998"
% ============================================================================
\section{Introducción}

\subsection{Propósito}

El propósito de este documento es especificar, de manera completa,
verificable y trazable, los requisitos funcionales y no funcionales del
sistema \textbf{LSA Salud}, un sistema de comunicación bidireccional entre
pacientes que usan Lengua de Señas Argentina (LSA) y personal de recepción
o enfermería, desplegado en el Sanatorio 9 de Julio, en entornos de
atención médica donde no hay un intérprete humano disponible. Sigue la
estructura recomendada por el estándar IEEE Std 830-1998. Los requisitos
(RF-01 a RF-06, RNF-01 a RNF-07) se desarrollan en sub-requisitos técnicos
trazables a lo largo de este documento (ver Sección~3.2 y
\texttt{IDS.md}).

Este documento está dirigido a:
\begin{itemize}
    \item El equipo de desarrollo del proyecto.
    \item Docentes y evaluadores académicos de la UTN--FRT.
    \item Profesionales de salud del Sanatorio 9 de Julio que operan el sistema.
    \item Futuros colaboradores o integradores del sistema.
\end{itemize}

\subsection{Ámbito del sistema}

\textbf{Nombre del sistema:} LSA Salud --- Sistema Inteligente de
Comunicación y Reconocimiento de LSA. \textbf{Cliente:} Sanatorio 9 de
Julio (San Miguel de Tucumán). El nombre del repositorio de código,
``Nexora'', es un alias técnico interno y no el nombre del proyecto
académico.

\subsubsection{Lo que el sistema hace}
\begin{itemize}
    \item Captura en tiempo real los movimientos corporales y los rasgos
    faciales del paciente mediante visión por computadora (MediaPipe:
    landmarks de manos, pose y rostro).
    \item Clasifica las señas capturadas con una red neuronal LSTM
    entrenada con datos propios.
    \item Reconstruye oraciones naturales en español a partir de las
    palabras sueltas reconocidas, usando plantillas médicas precalculadas
    y, como respaldo, un modelo de lenguaje local.
    \item Sintetiza las oraciones en voz con acento argentino.
    \item Provee un canal de comunicación bidireccional entre el paciente
    y el personal de recepción o enfermería, sincronizado entre dos estaciones de trabajo.
    \item Permite al personal de recepción o enfermería responder por
    texto o por dictado de voz.
    \item Ofrece un módulo de captura de datos para ampliar el vocabulario
    del modelo y un script de reentrenamiento.
    \item Solicita consentimiento explícito de cámara antes de iniciar la
    captura (RF-03.1), registra eventos de auditoría anónimos por
    operación relevante (RF-04.1) y permite reportar una traducción
    incorrecta o ambigua (RF-06.1). Ver Sección~3.2 y \texttt{IDS.md}.
\end{itemize}

\subsubsection{Lo que el sistema no hace}
\begin{itemize}
    \item No reemplaza al intérprete humano de LSA: actúa como un ajuste
    razonable complementario, sobre un vocabulario acotado.
    \item No almacena datos clínicos ni el historial de conversaciones: el
    estado de la conversación es volátil (en memoria del navegador, vía
    BroadcastChannel).
    \item No envía datos a servidores externos fuera de la máquina local
    (el modelo de lenguaje y la síntesis de voz corren localmente).
    \item No reconoce LSA como lengua continua y completa: reconoce el
    vocabulario clínico del dataset del proyecto, ampliable de forma
    continua mediante el Modo Captura (CU-06, CU-07).
\end{itemize}

\subsubsection{Beneficios, objetivos y metas}
\begin{itemize}
    \item Reducir la barrera de comunicación en centros de salud sin
    intérprete de LSA disponible.
    \item Preservar la privacidad de los datos de salud mediante
    procesamiento local.
    \item Demostrar la viabilidad de un flujo extremo a extremo de visión
    por computadora, procesamiento de lenguaje natural y síntesis de voz
    en hardware de consumo.
\end{itemize}

\subsection{Definiciones, acrónimos y abreviaturas}

\begin{longtable}{p{4cm} p{10cm}}
    \caption{Definiciones, acrónimos y abreviaturas}
    \label{tab:definiciones} \\
    \toprule
    \textbf{Término} & \textbf{Definición} \\
    \midrule
    \endfirsthead
    \toprule
    \textbf{Término} & \textbf{Definición} \\
    \midrule
    \endhead
    LSA & Lengua de Señas Argentina --- lengua natural con sintaxis propia. \\
    LSTM & Long Short-Term Memory --- red neuronal recurrente para secuencias temporales, usada para clasificar señas. \\
    LLM & Large Language Model --- modelo de lenguaje grande (Qwen2 1.5B, servido localmente por Ollama). \\
    TTS & Text-to-Speech --- síntesis de texto a voz (motor Piper, modelo ONNX). \\
    MediaPipe & Framework de Google para extracción de landmarks de manos, pose y rostro en tiempo real; pipeline de captura único de LSA Salud. \\
    Landmark & Punto anatómico detectado con coordenadas (x, y, z, visibilidad). \\
    Frame & Conjunto de 75 landmarks corporales (33 de pose, 21 por mano) más 52 blendshapes faciales, capturados en un instante. \\
    Blendshape & Coeficiente de expresión facial (0 a 1) provisto por MediaPipe; LSA Salud captura 52 por frame y los usa, junto con los landmarks de manos y pose, como características de reconocimiento: los blendshapes aportan los marcadores no manuales gramaticales de la LSA (preguntas, negación, intensidad) que los landmarks corporales por sí solos no capturan. \\
    Idle & Estado de reposo del paciente (manos abajo, sin movimiento significativo) que participa en el cierre de una oración. \\
    Plantilla médica & Frase médica predefinida asociada a una combinación conocida de clases de seña. \\
    Modo Captura & Modo de la Estación del Paciente para grabar nuevas muestras de entrenamiento. \\
    ERS / SRS & Especificación de Requisitos de Software / Software Requirements Specification. \\
    CU & Caso de Uso. \\
    RF / RNF & Requisito Funcional / Requisito No Funcional. \\
    REST & Representational State Transfer --- estilo arquitectónico para servicios web. \\
    WebSocket (WS) & Protocolo de comunicación bidireccional full-duplex sobre TCP. \\
    ONNX & Open Neural Network Exchange --- formato abierto para modelos de IA. \\
    BroadcastChannel API & API del navegador para comunicación entre pestañas del mismo origen, sin pasar por la red. \\
    p50 / p95 & Percentil 50 (mediana) y percentil 95 de una distribución de latencias. \\
    \bottomrule
\end{longtable}

\subsection{Referencias}

\begin{enumerate}
    \item IEEE Std 830-1998 --- IEEE Recommended Practice for Software Requirements Specifications.
    \item ISO/IEC/IEEE 42010 / UML 2.5 --- Object Management Group, \textit{Unified Modeling Language, Version 2.5.1}.
    \item Larman, C. \textit{UML y Patrones: Introducción al análisis y diseño orientado a objetos y al proceso unificado}. Pearson, 3.\textsuperscript{a} edición.
    \item Ley 26.378 --- Aprobación de la Convención sobre los Derechos de las Personas con Discapacidad (Argentina, 2008).
    \item Ley 27.044 --- Otorga jerarquía constitucional a la Convención (Argentina, 2014).
    \item Ley 25.326 --- Protección de Datos Personales (Argentina, 2000).
    \item Google MediaPipe --- \url{https://mediapipe.dev}
    \item PyTorch --- \url{https://pytorch.org}
    \item Ollama --- \url{https://ollama.ai}
    \item Piper TTS --- \url{https://github.com/rhasspy/piper}
    \item \texttt{README.md} y código fuente del repositorio ``Nexora'' (nombre interno del repositorio de código; \texttt{app.py}, \texttt{train\_lstm.py}, \texttt{src/\allowbreak model/\allowbreak lstm.py}, \texttt{web/}).
\end{enumerate}

\subsection{Visión general del documento}

El resto de este documento se organiza así:
\begin{itemize}
    \item \textbf{Sección 2 --- Descripción general:} contexto del
    producto, funciones generales, características de los usuarios,
    restricciones y suposiciones.
    \item \textbf{Sección 3 --- Requisitos específicos:} interfaces
    externas, requisitos funcionales, requisitos de rendimiento,
    restricciones de diseño y atributos del sistema.
    \item \textbf{Apéndices:} matriz de trazabilidad entre requisitos
    funcionales, casos de uso y requisitos no funcionales.
\end{itemize}

% ============================================================================
\section{Descripción general}

\subsection{Perspectiva del producto}

LSA Salud es un sistema \textbf{independiente} (\textit{standalone}) que
corre íntegramente en una PC local. No depende de servicios en la nube
para su funcionamiento (el modelo de lenguaje y la síntesis de voz corren
localmente); si el proceso de Ollama no está disponible, la traducción cae
en un modo degradado que igual entrega las palabras originales
(ver RF-01.5, extensión 3a.2a).

El sistema usa una \textbf{arquitectura de doble cara}: la Estación
del Paciente (cámara, predicción y chat) y la Estación de Recepción
(historial de chat y respuesta por texto o voz), sincronizadas mediante
BroadcastChannel. Como se detalla en la
Sección~2.5, estas dos caras son, a la vez, el diseño de la estación fija
de guardia/triaje exigida por los requisitos RNF-04 y RNF-07.

El sistema completo se despliega como \textbf{N estaciones fijas
independientes}, una por cada punto de atención del sanatorio (admisión,
recepción, triaje, guardia). Cada estación es una instancia local
autocontenida --- su propio navegador, servidor Next.js, backend FastAPI y
proceso Ollama --- sin servidor central ni datos compartidos entre
estaciones: cada una atiende su propia conversación de principio a fin y
la descarta al finalizar la atención (CU-12). Este diseño es consistente
con RNF-T-02 (el sistema no depende de conectividad externa) y con RNF-03
(sin almacenamiento persistente); agregar una estación nueva no requiere
coordinarla con las demás. El detalle de nodos y artefactos por estación
se muestra en la Vista Física (Entregable 06, Arquitectura).

\begin{table}[H]
    \centering
    \caption{Stack tecnológico de LSA Salud}
    \label{tab:stack}
    \small
    \begin{tabularx}{\textwidth}{>{\raggedright\arraybackslash}p{4cm} >{\raggedright\arraybackslash}p{4.4cm} X}
        \toprule
        \textbf{Componente} & \textbf{Tecnología} & \textbf{Función} \\
        \midrule
        Visión por computadora & MediaPipe & Extrae landmarks de manos, pose y rostro (blendshapes). \\
        Clasificación de señas & PyTorch (LSTM) & Clasifica ventanas de 45 frames en clases de seña. \\
        Reconstrucción gramatical & Plantillas médicas + Ollama (Qwen2 1.5B) & Convierte palabras sueltas en una oración natural. \\
        Síntesis de voz & Piper TTS (ONNX) & Verbaliza la oración con voz argentina. \\
        Reconocimiento de voz & Web Speech API & Permite al personal de recepción o enfermería dictar respuestas. \\
        Sincronización entre estaciones & BroadcastChannel API & Mensajes entre pestañas del mismo origen, sin red. \\
        Frontend & Next.js 14 + React + TypeScript & Interfaz web de las estaciones. \\
        Backend & FastAPI + Uvicorn & API REST y WebSocket para inferencia y traducción. \\
        \bottomrule
    \end{tabularx}
\end{table}

\subsection{Funciones del producto}

\begin{enumerate}
    \item \textbf{Traducción de LSA a texto:} captura de señas, clasificación por LSTM y reconstrucción de oraciones naturales (CU-01, CU-02).
    \item \textbf{Síntesis de voz:} conversión de la oración traducida a audio con voz argentina (CU-03).
    \item \textbf{Canal de comunicación bidireccional:} chat sincronizado entre las dos estaciones (CU-01, CU-04).
    \item \textbf{Respuesta por texto o por voz:} el personal de recepción o enfermería responde escribiendo o dictando (CU-04, CU-05).
    \item \textbf{Captura de datos de entrenamiento:} grabación de nuevas señas para ampliar el vocabulario (CU-06, CU-07).
    \item \textbf{Entrenamiento del modelo:} reentrenamiento del LSTM con los datos capturados (CU-08).
    \item \textbf{Consentimiento de cámara} (CU-09), \textbf{reporte de traducción incorrecta} (CU-10) y \textbf{registro de auditoría} (CU-11).
\end{enumerate}

\subsection{Características de los usuarios}

\begin{table}[H]
    \centering
    \caption{Tipos de usuario y perfiles}
    \label{tab:usuarios}
    \small
    \begin{tabularx}{\textwidth}{>{\raggedright\arraybackslash}p{3.6cm} >{\raggedright\arraybackslash}p{4.1cm} X}
        \toprule
        \textbf{Actor} & \textbf{Perfil} & \textbf{Interacción} \\
        \midrule
        Paciente (ACT-01) & Usuario de LSA, sin conocimiento técnico esperado. & Hace señas frente a la cámara; lee las respuestas del profesional en pantalla. \\
        Personal de Recepción o Enfermería (ACT-02) & Personal administrativo o de enfermería. & Lee las traducciones del paciente; responde por texto o por dictado de voz. \\
        Desarrollador de IA / Equipo de Proyecto (ACT-03) & Personal técnico del equipo del proyecto. & Usa el Modo Captura para grabar señas nuevas y ejecuta el reentrenamiento. \\
        \bottomrule
    \end{tabularx}
\end{table}

\subsection{Restricciones}

\begin{itemize}
    \item \textbf{Ejecución local obligatoria:} ningún dato de la conversación sale de la máquina; el único tráfico de red es hacia procesos locales (Ollama en \texttt{localhost:11434}).
    \item \textbf{Dependencia de Ollama:} el servicio debe estar corriendo localmente para la reconstrucción gramatical por LLM; si no está disponible, se degrada a texto sin reconstruir (RF-01.5, extensión 3a.2a).
    \item \textbf{Navegador:} Google Chrome es necesario para Web Speech API (dictado) y para un soporte confiable de MediaPipe.
    \item \textbf{Vocabulario clínico:} el modelo reconoce el vocabulario LSA clínico del dataset del proyecto (dataset ampliado desde LSA64, adaptado a síntomas, triaje, dolor y admisión/recepción/guardia, validado por intérpretes de LSA --- RNF-05/RNF-06), con señas como \texttt{cabeza}, \texttt{dolor}, \texttt{respirar} o \texttt{yo} como ejemplos; el vocabulario se amplía de forma continua mediante el Modo Captura y el reentrenamiento (CU-06, CU-07, CU-08).
    \item \textbf{Sin persistencia de conversaciones:} todo el historial de chat es volátil (memoria del navegador + BroadcastChannel).
    \item \textbf{Lenguajes:} Python 3.9+ en el backend; TypeScript/React en el frontend.
    \item \textbf{Hardware mínimo:} PC con cámara web; GPU opcional para acelerar la inferencia del LSTM (funciona sobre CPU).
\end{itemize}

\subsection{Suposiciones y dependencias}

\begin{itemize}
    \item Se asume que el sanatorio dispone, por cada punto de atención desplegado (admisión, recepción, triaje, guardia), de una PC con dos pantallas (o dos dispositivos) y una cámara web: las dos caras sincronizadas (Paciente y Recepción) implementan las dos caras de la estación fija de guardia/triaje exigida por RNF-04/RNF-07 (ver \texttt{IDS.md}, Sección~4.2); su ensamblaje físico sobre hardware comercial y su calibración se realizan durante la instalación de cada estación.
    \item Se asume que Ollama y el modelo \texttt{qwen2:1.5b} están instalados y descargados de antemano.
    \item Se asume que el modelo de voz Piper (\texttt{es\_AR-daniela-high.onnx}) está disponible en la carpeta \texttt{voices/}.
    \item Se asume que el paciente conoce LSA y puede realizar señas de forma clara frente a la cámara.
    \item Se asume conexión eléctrica estable; no se contempla operación a batería.
    \item Si cambia el sistema operativo (probado en Windows), puede requerirse revalidar la compatibilidad de MediaPipe y Piper.
\end{itemize}

% ============================================================================
\section{Requisitos específicos}

Esta sección detalla los requisitos con nivel suficiente para el diseño,
la implementación y la verificación del sistema. Cada requisito está
identificado de forma unívoca y es trazable al código fuente (ver
Apéndice~A).

\subsection{Interfaces externas}

\subsubsection{Interfaces de usuario}

\begin{table}[H]
    \centering
    \caption{Interfaces de usuario del sistema}
    \label{tab:interfaces-usuario}
    \begin{tabularx}{\textwidth}{l l X}
        \toprule
        \textbf{Interfaz} & \textbf{Ruta} & \textbf{Descripción} \\
        \midrule
        Landing page & \texttt{/} & Selector de estación (Paciente / Recepción). \\
        Diálogo de Consentimiento (P-05) & \texttt{/paciente} (modal previo) & Consentimiento explícito antes de activar la cámara. \\
        Estación del Paciente & \texttt{/paciente} & Cámara, predicción en tiempo real, chat y Modo Captura. \\
        Estación de Recepción & \texttt{/recepcion} & Historial de chat, respuesta por texto y por dictado de voz. \\
        \bottomrule
    \end{tabularx}
\end{table}

\subsubsection{Interfaces de hardware}

\begin{table}[H]
    \centering
    \caption{Interfaces de hardware}
    \label{tab:interfaces-hw}
    \begin{tabularx}{\textwidth}{l l X}
        \toprule
        \textbf{Hardware} & \textbf{API} & \textbf{Uso} \\
        \midrule
        Cámara web & \texttt{navigator.mediaDevices.getUserMedia()} & Captura de video para MediaPipe. \\
        Micrófono & Web Speech API (\texttt{SpeechRecognition}) & Dictado de voz del personal de recepción o enfermería. \\
        Parlantes & \texttt{HTMLAudioElement} & Reproducción de la voz sintetizada (Piper TTS). \\
        \bottomrule
    \end{tabularx}
\end{table}

\subsubsection{Interfaces de software}

\begin{table}[H]
    \centering
    \caption{Endpoints del backend FastAPI}
    \label{tab:endpoints}
    \begin{tabularx}{\textwidth}{l l X}
        \toprule
        \textbf{Método} & \textbf{Endpoint} & \textbf{Descripción} \\
        \midrule
        GET & \texttt{/} & Health check: estado, dispositivo de cómputo, cantidad de clases. \\
        GET & \texttt{/classes} & Lista de clases que reconoce el modelo cargado. \\
        GET & \texttt{/dataset\_stats} & Cantidad de muestras por etiqueta del dataset activo. \\
        POST & \texttt{/predict} & Predicción de una seña vía REST (frames en JSON). \\
        WS & \texttt{/ws/predict} & Predicción de señas vía WebSocket. \\
        POST & \texttt{/capture} & Guarda una muestra de entrenamiento (frames + etiqueta $\rightarrow$ \texttt{.npy}). \\
        POST & \texttt{/process\_sentence} & Traduce palabras sueltas a oración natural y genera el audio. \\
        \midrule
        POST & \texttt{/report} & Registra un reporte de traducción incorrecta o ambigua (CU-10, RF-06.1); invocado por \texttt{ChatBubble}. \\
        POST & \texttt{/audit} & Registra un evento de auditoría anónimo (CU-11, RF-04.1). Los tipos \texttt{consent} y \texttt{translation\_displayed} los envía el frontend (\texttt{PacientePage}, \texttt{useSentenceBuilder}); los tipos \texttt{prediction}, \texttt{sentence} y \texttt{report} los registra el propio backend al completar la operación correspondiente. \\
        GET & \texttt{/metrics?since=} & Consulta y agrega en SQL los eventos de auditoría de SQLite (\texttt{data/lsa\_salud.db}): conteos por tipo, consentimientos aceptados/declinados, fuente de traducción y latencias p50/p95/máx.\ por campo. \\
        GET & \texttt{/reports/summary?limit=} & Devuelve los mensajes y las palabras reconocidas más reportados (\texttt{limit} entre 1 y 100), agregados en SQL sobre \texttt{reported\_message}, \texttt{recognized\_word} y \texttt{translation\_report}, como insumo para priorizar el reentrenamiento (CU-08, RF-06.1). \\
        \bottomrule
    \end{tabularx}
\end{table}

También hay una dependencia de software externa: el servicio Ollama
(API REST local en el puerto 11434) que sirve el modelo de lenguaje
\texttt{qwen2:1.5b}.

\subsubsection{Interfaces de comunicación}

\begin{itemize}
    \item \textbf{BroadcastChannel:} canal \texttt{nexora-chat} (nombre interno heredado del repositorio de código) para sincronizar mensajes entre pestañas del mismo origen. Eventos: \texttt{MESSAGE} (nuevo mensaje) y \texttt{CLEAR} (limpiar historial), con deduplicación por identificador de mensaje.
    \item \textbf{WebSocket:} conexión persistente entre el frontend y \texttt{/ws/predict}, con reconexión automática cada 2 segundos ante corte.
    \item \textbf{HTTP REST:} comunicación solicitud-respuesta para traducción de oraciones, captura de datos y consulta de clases/estadísticas.
\end{itemize}

\subsection{Requisitos funcionales}

Los seis requisitos funcionales base (RF-01 a RF-06) se desarrollan aquí
como sub-requisitos técnicos trazables, numerados \texttt{RF-<base>.<n>}
(ver \texttt{IDS.md}, Sección~3, para el texto literal de cada RF y la
tabla de trazabilidad completa). Los requisitos sin RF base --- decisión
de diseño del equipo para CU-06/07/08 --- usan el prefijo \texttt{RF-T-}.
CU-06, CU-07 y CU-08 son, en conjunto, las herramientas de administración
del modelo y del dataset con las que cuenta el Equipo de Proyecto (Modo
Captura, gestión de vocabulario y reentrenamiento); el sistema no incluye
una pantalla de administración adicional fuera de estos tres casos de uso.
Cada requisito indica su prioridad (\textit{Esencial}: imprescindible para
el flujo principal; \textit{Condicional}: mejora el flujo pero el sistema
es utilizable sin él; \textit{Deseable}: conveniencia menor) y los casos
de uso relacionados (ver Entregable de Modelos de Casos de Uso). Al final
de la sección se detallan RF-03.1, RF-04.1 y RF-06.1.

\subsubsection{RF-01.1: Captura de landmarks en tiempo real (RF-01)}
\begin{table}[H]
\centering
\begin{tabularx}{\textwidth}{l X}
\toprule
\textbf{ID} & RF-01.1 \\
\textbf{Prioridad} & Esencial \\
\textbf{Descripción} & El sistema debe extraer, en tiempo real, los landmarks corporales del paciente desde la cámara web. \\
\textbf{Entrada} & Flujo de video de la cámara ($\sim$30 fps). \\
\textbf{Proceso} & MediaPipe procesa cada frame elegible y extrae landmarks de pose (33) y de ambas manos (21 cada una), normalizados respecto de la nariz, junto con 52 blendshapes faciales que capturan los marcadores no manuales gramaticales de la LSA. \\
\textbf{Salida} & Un \textit{frame} de 75 landmarks corporales (300 valores) más 52 blendshapes faciales, por cada muestreo. \\
\textbf{CU relacionados} & CU-01. \\
\bottomrule
\end{tabularx}
\end{table}

\subsubsection{RF-01.2: Envío de frames para predicción (RF-01)}
\begin{table}[H]
\centering
\begin{tabularx}{\textwidth}{l X}
\toprule
\textbf{ID} & RF-01.2 \\
\textbf{Prioridad} & Esencial \\
\textbf{Descripción} & El frontend debe acumular frames en un buffer y enviarlos periódicamente al backend para su clasificación. \\
\textbf{Entrada} & Buffer de frames (hasta 45, con al menos 20 acumulados antes del primer envío). \\
\textbf{Proceso} & Cada 8 frames nuevos, el buffer se serializa y se envía por WebSocket al backend. \\
\textbf{Salida} & Mensaje WebSocket entrante con la predicción (ver RF-01.3). \\
\textbf{CU relacionados} & CU-01. \\
\bottomrule
\end{tabularx}
\end{table}

\subsubsection{RF-01.3: Clasificación de la seña mediante LSTM (RF-01; realiza también RF-05)}
\begin{table}[H]
\centering
\begin{tabularx}{\textwidth}{l X}
\toprule
\textbf{ID} & RF-01.3 \\
\textbf{Prioridad} & Esencial \\
\textbf{Descripción} & El backend debe clasificar la ventana de frames recibida en una de las clases de seña conocidas, con un nivel de confianza. El umbral de confianza (ver ``Umbral'' más abajo) es, a la vez, el mecanismo que realiza el requisito RF-05 (``detectar cuando una seña no pueda interpretarse con confianza y evitar mostrar una traducción incorrecta''). \\
\textbf{Entrada} & Tensor de dimensiones $(1, 45, 300)$, recortado o relleno con ceros según haga falta. \\
\textbf{Proceso} & Red LSTM (2 capas, 128 unidades ocultas) seguida de una capa totalmente conectada y \textit{softmax}; se calcula la clase de mayor probabilidad y el top-3. \\
\textbf{Salida} & Seña predicha, confianza (\%) y top-3 de predicciones. \\
\textbf{Umbral} & $\geq$70\% de confianza para que la palabra se registre en la secuencia (ver RF-01.4). \\
\textbf{CU relacionados} & CU-01, CU-02. \\
\bottomrule
\end{tabularx}
\end{table}

\subsubsection{RF-01.4: Construcción de la secuencia y cierre de oración (RF-01; realiza también RF-05)}
\begin{table}[H]
\centering
\begin{tabularx}{\textwidth}{l X}
\toprule
\textbf{ID} & RF-01.4 \\
\textbf{Prioridad} & Esencial \\
\textbf{Descripción} & El sistema debe acumular las palabras reconocidas en una secuencia y detectar cuándo el paciente terminó de expresar su idea, descartando las señas de baja confianza o de reposo (mecanismo que también realiza RF-05). \\
\textbf{Entrada} & Predicciones sucesivas del LSTM y magnitud de movimiento entre frames consecutivos. \\
\textbf{Proceso} & Se agrega una palabra a la secuencia cuando su confianza es $\geq$70\% y no es la clase de reposo (\texttt{idle}) ni domina el top-3 con más de 50\%. La secuencia se cierra cuando la magnitud media de movimiento de los últimos 6 frames es menor a 0.05 y pasaron al menos 800 ms desde la última predicción distinta de reposo. \\
\textbf{Salida} & Secuencia de palabras cerrada, lista para traducirse (dispara CU-02). \\
\textbf{CU relacionados} & CU-01. \\
\bottomrule
\end{tabularx}
\end{table}

\subsubsection{RF-01.5: Traducción de la secuencia a oración natural (RF-01)}
\begin{table}[H]
\centering
\begin{tabularx}{\textwidth}{l X}
\toprule
\textbf{ID} & RF-01.5 \\
\textbf{Prioridad} & Esencial \\
\textbf{Descripción} & El sistema debe convertir la secuencia de palabras sueltas en una oración en español natural, priorizando respuestas instantáneas. \\
\textbf{Entrada} & Cadena de palabras sueltas (por ejemplo, ``yo dolor cabeza''). \\
\textbf{Proceso} & Se normaliza la secuencia y se busca, en orden: coincidencia exacta en el diccionario de plantillas médicas; luego en la caché de traducciones recientes (hasta 200 entradas); si no hay coincidencia, se delega en el modelo de lenguaje local (tiempo de espera de 8 segundos), con instrucciones explícitas de no inventar síntomas y de respetar la persona gramatical. \\
\textbf{Salida} & Oración natural (por ejemplo, ``Me duele la cabeza.''), y el origen de la traducción (plantilla, caché o modelo de lenguaje). \\
\textbf{CU relacionados} & CU-02. \\
\bottomrule
\end{tabularx}
\end{table}

\subsubsection{RF-01.6: Síntesis de voz (RF-01)}
\begin{table}[H]
\centering
\begin{tabularx}{\textwidth}{l X}
\toprule
\textbf{ID} & RF-01.6 \\
\textbf{Prioridad} & Esencial \\
\textbf{Descripción} & La oración traducida debe sintetizarse como audio, con voz en español argentino. \\
\textbf{Entrada} & Oración en español (texto). \\
\textbf{Proceso} & Si el audio ya fue generado antes para la misma oración, se reutiliza desde la caché (hasta 50 entradas); en caso contrario, se sintetiza con Piper TTS y el resultado se codifica en base64. \\
\textbf{Salida} & Audio WAV codificado en base64, reproducido en el navegador tras una ventana de exhibición del texto (1000 ms por defecto). \\
\textbf{CU relacionados} & CU-03. \\
\bottomrule
\end{tabularx}
\end{table}

\subsubsection{RF-01.7: Sincronización de la conversación entre estaciones (RF-01; soporta también RF-02)}
\begin{table}[H]
\centering
\begin{tabularx}{\textwidth}{l X}
\toprule
\textbf{ID} & RF-01.7 \\
\textbf{Prioridad} & Esencial \\
\textbf{Descripción} & Los mensajes del paciente y del personal de recepción o enfermería deben sincronizarse instantáneamente entre las dos estaciones. Es el canal que entrega el mensaje del profesional en un formato accesible para el paciente sordo (RF-02). \\
\textbf{Entrada} & Mensaje nuevo (del paciente, generado por CU-02/CU-03, o del personal, generado por CU-04). \\
\textbf{Proceso} & El mensaje se agrega al historial local y se publica por BroadcastChannel; cada estación deduplica por identificador antes de mostrarlo. \\
\textbf{Salida} & Historial de conversación actualizado en ambas estaciones. \\
\textbf{CU relacionados} & CU-01, CU-04. \\
\bottomrule
\end{tabularx}
\end{table}

\subsubsection{RF-01.8: Verificación de fidelidad clínica de la traducción (RF-01; realiza también RF-05)}
\begin{table}[H]
\centering
\begin{tabularx}{\textwidth}{l X}
\toprule
\textbf{ID} & RF-01.8 \\
\textbf{Prioridad} & Esencial. \\
\textbf{Descripción} & El módulo de traducción debe validar que la oración generada por el modelo de lenguaje no incorpore síntomas, medicamentos ni términos clínicos ausentes de la secuencia de señas reconocida, antes de mostrarla o sintetizarla. \\
\textbf{Entrada} & Secuencia de palabras reconocida y oración candidata generada por el modelo de lenguaje (RF-01.5). \\
\textbf{Proceso} & El prompt enviado al modelo de lenguaje local instruye explícitamente no inventar síntomas y respetar la persona gramatical (ver RF-01.5); el sistema verifica que los términos clínicos de la oración generada estén contenidos en la secuencia reconocida antes de aceptarla. \\
\textbf{Salida} & Oración validada, o texto original de la secuencia si la validación falla. \\
\textbf{CU relacionados} & CU-02. \\
\bottomrule
\end{tabularx}
\end{table}

\subsubsection{RF-02.1: Respuesta del personal por texto (RF-02)}
\begin{table}[H]
\centering
\begin{tabularx}{\textwidth}{l X}
\toprule
\textbf{ID} & RF-02.1 \\
\textbf{Prioridad} & Esencial \\
\textbf{Descripción} & El personal de recepción o enfermería debe poder escribir una respuesta y enviarla a la conversación compartida. \\
\textbf{Entrada} & Texto ingresado por el personal. \\
\textbf{Proceso} & Se valida que el texto no esté vacío antes de habilitar el envío. \\
\textbf{Salida} & Mensaje agregado a la conversación (ver RF-01.7). \\
\textbf{CU relacionados} & CU-04. \\
\bottomrule
\end{tabularx}
\end{table}

\subsubsection{RF-02.2: Dictado por voz del personal (RF-02)}
\begin{table}[H]
\centering
\begin{tabularx}{\textwidth}{l X}
\toprule
\textbf{ID} & RF-02.2 \\
\textbf{Prioridad} & Condicional \\
\textbf{Descripción} & El personal de recepción o enfermería debe poder dictar su respuesta por voz en lugar de escribirla. \\
\textbf{Entrada} & Audio del micrófono. \\
\textbf{Proceso} & Se usa la Web Speech API del navegador, configurada en español argentino, con resultados parciales y finales; el texto dictado se acumula sobre el texto ya escrito. \\
\textbf{Salida} & Texto transcripto, listo para enviarse (RF-02.1). \\
\textbf{Restricción} & Requiere un navegador compatible (Google Chrome). \\
\textbf{CU relacionados} & CU-05. \\
\bottomrule
\end{tabularx}
\end{table}

\subsubsection{RF-T-01: Selección o creación de etiqueta de seña (decisión de diseño del equipo)}
\begin{table}[H]
\centering
\begin{tabularx}{\textwidth}{l X}
\toprule
\textbf{ID} & RF-T-01 \\
\textbf{Prioridad} & Condicional \\
\textbf{Descripción} & Antes de grabar una muestra, el desarrollador debe poder elegir una etiqueta existente o crear una nueva. \\
\textbf{Entrada} & Nombre de etiqueta (existente o nueva). \\
\textbf{Proceso} & Si la etiqueta es nueva, se valida contra un patrón (letras minúsculas, dígitos y guion bajo, de 1 a 32 caracteres) antes de aceptarla; la carpeta de datos se crea recién al guardar la primera muestra. \\
\textbf{Salida} & Etiqueta activa para la siguiente grabación. \\
\textbf{CU relacionados} & CU-07. \\
\bottomrule
\end{tabularx}
\end{table}

\subsubsection{RF-T-02: Consulta de clases y estadísticas del dataset (decisión de diseño del equipo)}
\begin{table}[H]
\centering
\begin{tabularx}{\textwidth}{l X}
\toprule
\textbf{ID} & RF-T-02 \\
\textbf{Prioridad} & Condicional \\
\textbf{Descripción} & El sistema debe exponer las clases que reconoce el modelo y la cantidad de muestras grabadas por etiqueta. \\
\textbf{Entrada} & Solicitud de estadísticas del dataset activo. \\
\textbf{Proceso} & El backend lista las clases del modelo cargado y cuenta los archivos \texttt{.npy} por carpeta de etiqueta en el árbol de datos del dataset. \\
\textbf{Salida} & Lista de clases y mapa etiqueta $\rightarrow$ cantidad de muestras. \\
\textbf{CU relacionados} & CU-07. \\
\bottomrule
\end{tabularx}
\end{table}

\subsubsection{RF-T-03: Grabación de muestras de entrenamiento (decisión de diseño del equipo)}
\begin{table}[H]
\centering
\begin{tabularx}{\textwidth}{l X}
\toprule
\textbf{ID} & RF-T-03 \\
\textbf{Prioridad} & Condicional \\
\textbf{Descripción} & El sistema debe permitir grabar una o varias muestras de una seña, con cuenta regresiva y repeticiones configurables. \\
\textbf{Entrada} & Etiqueta activa, cantidad de repeticiones. \\
\textbf{Proceso} & Por cada repetición: cuenta regresiva de 3 segundos, grabación de 45 frames consecutivos, envío al backend. \\
\textbf{Salida} & Confirmación de guardado y conteo de muestras actualizado por etiqueta. \\
\textbf{CU relacionados} & CU-06. \\
\bottomrule
\end{tabularx}
\end{table}

\subsubsection{RF-T-04: Persistencia de muestras de entrenamiento (decisión de diseño del equipo)}
\begin{table}[H]
\centering
\begin{tabularx}{\textwidth}{l X}
\toprule
\textbf{ID} & RF-T-04 \\
\textbf{Prioridad} & Condicional \\
\textbf{Descripción} & Cada muestra grabada debe guardarse en disco, con los landmarks de manos y pose junto con los blendshapes faciales en un único árbol de captura. \\
\textbf{Entrada} & Frames grabados, con sus landmarks y blendshapes faciales. \\
\textbf{Proceso} & Se valida que los frames sean numéricamente finitos y que haya exactamente 52 blendshapes por frame; se guarda como archivo \texttt{.npy} bajo \texttt{dataset/custom/<etiqueta>/}. \\
\textbf{Salida} & Archivo \texttt{.npy} persistido con nombre de marca de tiempo. \\
\textbf{CU relacionados} & CU-06. \\
\bottomrule
\end{tabularx}
\end{table}

\subsubsection{RF-T-05: Reentrenamiento del modelo LSTM (decisión de diseño del equipo)}
\begin{table}[H]
\centering
\begin{tabularx}{\textwidth}{l X}
\toprule
\textbf{ID} & RF-T-05 \\
\textbf{Prioridad} & Condicional \\
\textbf{Descripción} & El sistema debe permitir reentrenar el modelo LSTM con las muestras capturadas, fuera de línea. \\
\textbf{Entrada} & Conjunto de archivos \texttt{.npy} agrupados por etiqueta en \texttt{dataset/custom}. \\
\textbf{Proceso} & Descubrimiento de clases, separación 80/20 entrenamiento/validación (semilla fija), aumento de datos en entrenamiento (recorte aleatorio, ruido gaussiano, escalado), 50 rondas de entrenamiento, conservación del mejor modelo por exactitud de validación. \\
\textbf{Salida} & Archivo de pesos del modelo y archivo de clases actualizados. \\
\textbf{CU relacionados} & CU-08. \\
\bottomrule
\end{tabularx}
\end{table}

\subsubsection{RF-T-06: Selección de estación de trabajo (decisión de diseño del equipo)}
\begin{table}[H]
\centering
\begin{tabularx}{\textwidth}{l X}
\toprule
\textbf{ID} & RF-T-06 \\
\textbf{Prioridad} & Deseable \\
\textbf{Descripción} & Desde la página de inicio, el usuario debe poder abrir la estación que corresponda a su rol. \\
\textbf{Entrada} & Selección del usuario. \\
\textbf{Proceso} & Navegación a \texttt{/paciente} o \texttt{/recepcion}. \\
\textbf{Salida} & Estación correspondiente abierta. \\
\textbf{CU relacionados} & Precondición de CU-01 y de CU-04. \\
\bottomrule
\end{tabularx}
\end{table}

\subsubsection{RF-03.1: Consentimiento explícito de cámara (RF-03)}
\begin{table}[H]
\centering
\begin{tabularx}{\textwidth}{l X}
\toprule
\textbf{ID} & RF-03.1 \\
\textbf{Prioridad} & Esencial (bloqueante: sin consentimiento, no hay CU-01). \\
\textbf{Descripción} & Antes de activar la cámara, el sistema despliega un diálogo modal de consentimiento explícito (P-05) sobre \texttt{PacientePage} y mantiene el resto de la página inerte (\texttt{inert}) hasta que el paciente decide. El estado de consentimiento (\texttt{pending} / \texttt{accepted} / \texttt{declined}) vive en memoria del componente, sin persistencia entre sesiones. \\
\textbf{Entrada} & Decisión del paciente (acepta / no acepta), capturada por el diálogo. \\
\textbf{Proceso} & \texttt{PacientePage} renderiza el diálogo y solo pasa \texttt{enabled: true} al hook de captura de landmarks (\texttt{useMediaPipe} / \texttt{useMediaPipeTasks}) cuando \texttt{consentState === "accepted"}; ese \texttt{enabled} gatea el efecto que invoca \texttt{getUserMedia()}, por lo que la cámara nunca se activa sin consentimiento previo. Cada decisión se registra como un evento de auditoría anónimo (\texttt{type: "consent"}, campo \texttt{consent}: \texttt{accepted}/\texttt{declined}) vía \texttt{POST /audit}. Si el paciente no acepta, puede reconsiderar más tarde desde el mismo diálogo. \\
\textbf{Salida} & Consentimiento registrado en la base de datos SQLite (\texttt{data/lsa\_salud.db}) y cámara activada, o cámara inactiva con opción de reconsiderar. \\
\textbf{CU relacionados} & CU-09 (precondición de CU-01). \\
\bottomrule
\end{tabularx}
\end{table}

\subsubsection{RF-04.1: Registro de eventos de auditoría (RF-04)}
\begin{table}[H]
\centering
\begin{tabularx}{\textwidth}{l X}
\toprule
\textbf{ID} & RF-04.1 \\
\textbf{Prioridad} & Condicional (no bloquea el flujo de comunicación; es una subfunción de auditoría). \\
\textbf{Descripción} & El sistema registra, de forma automática y anónima, un evento de auditoría por cada operación de sistema relevante (predicción, traducción/síntesis, reporte, consentimiento, traducción mostrada, respuesta enviada, fin de atención), sin datos identificatorios ni imágenes, video o texto de la conversación, para habilitar métricas operativas de uso . \\
\textbf{Entrada} & Finalización de una operación de sistema (SO-01 a SO-05, ver Entregable 08) o un evento anónimo enviado por el frontend (\texttt{consent}, \texttt{translation\_displayed}). \\
\textbf{Proceso} & El backend (función envoltura \texttt{\_log\_audit} en \texttt{app.py} que delega en \texttt{storage.log\_audit\_event}) inserta el registro estructurado en la tabla \texttt{audit\_event} de SQLite (\texttt{data/lsa\_salud.db}). En la ruta WebSocket (\texttt{/ws/predict}) la escritura se despacha con \texttt{asyncio.to\_thread} para no bloquear el bucle de eventos. \texttt{GET /metrics} ejecuta agregaciones SQL sobre la tabla (conteos por tipo, consentimientos, fuente de traducción, p50/p95/máx.\ de latencia). \\
\textbf{Salida} & Evento de auditoría anónimo persistido en SQLite (\texttt{data/lsa\_salud.db}). \\
\textbf{CU relacionados} & CU-11 (incluido por CU-01 y CU-04). \\
\bottomrule
\end{tabularx}
\end{table}

\subsubsection{RF-06.1: Reporte de traducción incorrecta o ambigua (RF-06)}
\begin{table}[H]
\centering
\begin{tabularx}{\textwidth}{l X}
\toprule
\textbf{ID} & RF-06.1 \\
\textbf{Prioridad} & Condicional. \\
\textbf{Descripción} & El paciente o el personal de recepción o enfermería pueden marcar un mensaje de la conversación (traducido del paciente) como traducción incorrecta o ambigua, para registrar el caso y priorizar el reentrenamiento (CU-08). \\
\textbf{Entrada} & Click en ``Reportar traducción'' sobre un mensaje del paciente, en \texttt{ChatBubble}. \\
\textbf{Proceso} & \texttt{ChatBubble} envía \texttt{POST /report} con el texto del mensaje, la estación (\texttt{paciente}/\texttt{recepcion}, derivada de la ruta con \texttt{usePathname}) y, si están disponibles, las palabras crudas reconocidas y el ID del mensaje. El backend valida el payload (\texttt{validate\_report\_payload}) y delega en \texttt{storage.save\_report} la inserción en las tablas \texttt{reported\_message}, \texttt{recognized\_word} y \texttt{translation\_report} de SQLite (\texttt{data/lsa\_salud.db}), además de registrar un evento de auditoría \texttt{report}. El frontend emite el evento \texttt{REPORTED} vía \texttt{BroadcastChannel} para actualizar y tachar el mensaje en ambas caras de la estación. \\
\textbf{Salida} & Reporte y mensaje persistidos en SQLite, notificación sincronizada a la otra cara y burbuja en estado permanente ``Traducción reportada'' con aviso operativo. \\
\textbf{CU relacionados} & CU-10 (extiende a CU-01 y CU-04). \\
\bottomrule
\end{tabularx}
\end{table}

\subsection{Requisitos de rendimiento (RNF-01.1, sub-requisito de RNF-01)}

Los siguientes valores son metas de diseño, expresadas en p50/p95, y
trazables al requisito RNF-01 (latencia percibida menor a 1 segundo). El
sistema mide cada etapa de
forma continua: cada predicción y cada oración procesada registra su
latencia por etapa (\texttt{lstm\_ms}, \texttt{translation\_ms},
\texttt{tts\_ms}, \texttt{total\_ms}) como evento de auditoría en
SQLite (\texttt{data/lsa\_salud.db}), y \texttt{GET /metrics} agrega esos campos en
p50/p95/máx.\ bajo demanda; el frontend también reporta \texttt{e2e\_ms}
(latencia percibida, desde el cierre de la oración hasta que se muestra
en pantalla) vía \texttt{POST /audit} (\texttt{translation\_displayed}).
Este criterio de verificación de RNF-01 se comprueba mediante esta
medición continua de p50/p95 sobre las métricas operativas del sistema.

\begin{table}[H]
    \centering
    \caption{Requisitos de rendimiento (RNF-01.1)}
    \label{tab:rendimiento}
    \small
    \begin{tabularx}{\textwidth}{>{\raggedright\arraybackslash}p{5.2cm} >{\raggedright\arraybackslash}p{3.6cm} X}
        \toprule
        \textbf{Métrica} & \textbf{Valor objetivo} & \textbf{Verificación} \\
        \midrule
        Latencia extremo a extremo (captura $\rightarrow$ audio listo), ruta con plantilla & p50 $\leq$ 600 ms & Medición continua vía \texttt{GET /metrics} \\
        Latencia extremo a extremo, caso general & p95 $\leq$ 1 s & Medición continua vía \texttt{GET /metrics} \\
        Inferencia LSTM por ventana (\texttt{lstm\_ms}) & sin regresión de p95 mayor al 10\% entre versiones del modelo & Medición continua vía \texttt{GET /metrics}, comparada entre versiones \\
        Captura de frames & $\sim$30 fps (según cámara) & Verificada por hardware de cámara estándar (RNF-04) \\
        Usuarios simultáneos & 1 par paciente--personal por instancia de estación & Verificada por el diseño de estación independiente (ver Sección~2.1) \\
        \bottomrule
    \end{tabularx}
\end{table}

\textbf{RNF-02.1 (sub-requisito de RNF-02, precisión):} la
exactitud de validación del modelo LSTM (\texttt{ModeloDeReconocimiento.exactitudDeValidación},
ver Entregable 01) se calcula durante \texttt{train\_lstm.py} sobre el 20\%
de validación. Además, el script \texttt{eval\_model.py <directorio>}
evalúa el modelo cargado por \texttt{app.py} sobre un directorio de
muestras \texttt{.npy}, reportando accuracy global, precisión/recall por
clase, matriz de confusión y cantidad de muestras bajo el umbral de
confianza; advierte si el directorio evaluado es el mismo usado para
entrenar (\texttt{dataset/custom}), ya que en ese caso la accuracy no mide
generalización real. El criterio de verificación de RNF-02 (accuracy
$\geq$ 90\% en entornos controlados y $\geq$ 85\% en operación real) se
comprueba corriendo \texttt{eval\_model.py} sobre una sesión de grabación
independiente del conjunto de entrenamiento.

\subsection{Restricciones de diseño}

\begin{itemize}
    \item \textbf{Privacidad por diseño:} todo el procesamiento sensible ocurre en la máquina local; el único proceso externo invocado es Ollama, también local.
    \item \textbf{Backend de un solo archivo:} \texttt{app.py} concentra los endpoints REST y WebSocket, pensado para simplicidad de despliegue en un entorno de consultorio.
    \item \textbf{Modelos pre-entrenados como dependencia:} los pesos del LSTM, el modelo de voz Piper y el modelo de lenguaje de Ollama deben estar descargados de antemano; el sistema no los descarga en tiempo de ejecución.
    \item \textbf{Protocolos por tipo de operación:} WebSocket para predicción de baja latencia, REST para operaciones puntuales (traducción, captura, consulta de clases), BroadcastChannel para sincronización entre estaciones.
    \item \textbf{Persistencia híbrida local:} SQLite embebido (\texttt{data/lsa\_salud.db}) para reportes y auditoría; muestras de entrenamiento como archivos \texttt{.npy} en el sistema de archivos; la conversación médica ordinaria es volátil y no se persiste.
\end{itemize}

\subsection{Atributos del sistema}

Los requisitos de rendimiento se documentan en la sección 3.3 (RNF-01.1).
Esta sección cubre los restantes atributos de calidad: los primeros tres
(RNF-03.1, RNF-04.1/RNF-07.1, RNF-05.1/RNF-06.1) son sub-requisitos
trazables a la matriz de requisitos; los seis siguientes (RNF-T-01 a
RNF-T-06) son requisitos técnicos definidos por el equipo, sin RNF base.

\subsubsection{Seguridad y privacidad (RNF-03.1, sub-requisito de RNF-03)}
\begin{itemize}
    \item \textbf{Privacidad de datos de salud:} todo el procesamiento y almacenamiento ocurre en la máquina local; no hay persistencia de conversaciones (estado en memoria más BroadcastChannel); al finalizar cada atención (CU-12), el sistema borra explícitamente la conversación en las dos caras de la estación.
    \item \textbf{Configuración de CORS:} el backend admite conexiones únicamente desde los orígenes configurados para la red local del sanatorio (ver RNF-T-06, Sección~3.4).
    \item \textbf{Acceso a cámara y micrófono:} requiere autorización explícita del navegador, otorgada por el usuario, y consentimiento explícito previo del paciente (CU-09, RF-03.1).
    \item \textbf{Eventos de auditoría (RF-04.1) y reportes de traducción (RF-06.1):} se registran en SQLite sin datos identificatorios ni imágenes o video; el texto del mensaje se guarda únicamente para los mensajes explícitamente reportados, como insumo técnico para priorizar el reentrenamiento del modelo.
\end{itemize}

\subsubsection{Compatibilidad de hardware y estación fija (RNF-04.1 / RNF-07.1, sub-requisitos de RNF-04 y RNF-07)}
\begin{itemize}
    \item Las dos caras sincronizadas de la estación (Paciente y Recepción), sobre el mismo navegador y canal BroadcastChannel, con doble monitor, implementan las dos caras de la estación fija de guardia/triaje que exigen RNF-04 y RNF-07 (cámara, pantalla doble cara, micrófono, altavoces, componentes comerciales estándar).
    \item \textbf{Calibración e instalación:} durante la instalación de cada estación, se realiza el ensamblaje físico de las dos pantallas y la cámara sobre un mostrador o box existente, con hardware comercial, y su calibración (encuadre de cámara, alineación de pantallas), antes de que la estación entre en servicio. Este procedimiento operativo es una restricción de instalación del sistema completo, no una pantalla de software.
\end{itemize}

\subsubsection{Especialización y validación del dataset (RNF-05.1 / RNF-06.1, sub-requisitos de RNF-05 y RNF-06)}
\begin{itemize}
    \item El Modo Captura y el árbol de captura del dataset (\texttt{dataset/custom}) permiten ampliar de forma continua el vocabulario clínico LSA (síntomas, triaje, dolor, admisión/recepción/guardia) a partir de un conjunto inicial de referencia (RNF-05.1).
    \item La validación lingüística del dataset por intérpretes de LSA es un procedimiento operativo del proyecto, realizado antes de incorporar cada ampliación del vocabulario al modelo en producción, y documentado mediante un acta de validación firmada por los intérpretes participantes (RNF-06.1).
\end{itemize}

\subsubsection{RNF-T-01: Fiabilidad (requisito técnico del equipo)}
\begin{itemize}
    \item El WebSocket de predicción se reconecta automáticamente cada 2 segundos ante un cierre de conexión.
    \item Si el modelo de lenguaje local falla o no responde a tiempo (8 s), se usa el texto original de la secuencia sin interrumpir el flujo.
    \item Si el motor Piper TTS no está cargado, el mensaje se entrega sin audio.
    \item Un indicador de solicitud en curso evita enviar frames al WebSocket mientras una predicción anterior sigue pendiente.
\end{itemize}

\subsubsection{RNF-T-02: Disponibilidad (requisito técnico del equipo)}
\begin{itemize}
    \item El sistema no depende de conectividad a internet para el flujo crítico (captura, clasificación LSTM, síntesis de voz): todos esos componentes corren localmente.
    \item La disponibilidad de la traducción por modelo de lenguaje depende de que el proceso local de Ollama esté activo; su ausencia degrada la traducción pero no interrumpe el sistema (ver RNF-T-01).
\end{itemize}

\subsubsection{RNF-T-03: Mantenibilidad (requisito técnico del equipo)}
\begin{itemize}
    \item Separación clara entre frontend (Next.js) y backend (FastAPI).
    \item Responsabilidades divididas en hooks de React independientes: captura de cámara, WebSocket de predicción, construcción de oraciones, canal de chat y reconocimiento de voz.
    \item El modelo LSTM está desacoplado del backend en un módulo propio (\texttt{src/\allowbreak model/\allowbreak lstm.py}).
    \item Las plantillas médicas están centralizadas en una estructura de datos del backend, fácil de extender.
\end{itemize}

\subsubsection{RNF-T-04: Portabilidad (requisito técnico del equipo)}
\begin{itemize}
    \item Probado en Windows; el stack (Next.js, FastAPI/Uvicorn, PyTorch, ONNX) es multiplataforma en principio.
    \item Requiere Node.js 18 o superior y Python 3.9 o superior.
    \item Navegador soportado: Google Chrome, requerido por la Web Speech API que usa RF-02.2 (dictado); el resto del pipeline (MediaPipe, WebSocket, TTS local) no depende de APIs propietarias de un navegador en particular.
\end{itemize}

\subsubsection{RNF-T-05: Usabilidad y accesibilidad (requisito técnico del equipo)}
\begin{itemize}
    \item Conversación fluida, sin botones de confirmación manual: la ventana de exhibición del texto antes de reproducir el audio actúa como confirmación automática (excepción deliberada: el consentimiento de cámara de P-05/CU-09 y el reporte de traducción de CU-10 sí requieren una acción explícita, ver Entregable 04).
    \item Interfaz en español, con indicadores visuales de estado (``EN VIVO'', ``GRABANDO'', cuenta regresiva, barras de confianza en tiempo real).
    \item Arquitectura de doble cara pensada para separar visualmente el rol del paciente y el del profesional de salud.
\end{itemize}

\subsubsection{RNF-T-06: Seguridad de red (requisito técnico del equipo)}
\begin{itemize}
    \item El backend admite conexiones únicamente desde los orígenes configurados para la red local del sanatorio (restricción de CORS); no expone el servicio a redes externas.
\end{itemize}

\subsection{Otros requisitos}

\subsubsection{Escalabilidad del vocabulario}
\begin{itemize}
    \item El Modo Captura permite ampliar el vocabulario de forma incremental, sin requerir conocimientos técnicos avanzados (CU-06, CU-07).
    \item El reentrenamiento es offline mediante \texttt{train\_lstm.py} (CU-08).
\end{itemize}

% ============================================================================
\section{Apéndices}

\subsection{Apéndice A: Matriz de trazabilidad RF -- CU -- RNF}

\small
\begin{longtable}{>{\raggedright\arraybackslash}p{3.4cm} >{\raggedright\arraybackslash}p{5.6cm} >{\raggedright\arraybackslash}p{5.6cm}}
    \caption{Matriz de trazabilidad Requisitos Funcionales -- Casos de Uso -- Requisitos No Funcionales}
    \label{tab:trazabilidad} \\
    \toprule
    \textbf{Requisito funcional} & \textbf{Casos de uso relacionados} & \textbf{Requisitos no funcionales relacionados} \\
    \midrule
    \endfirsthead
    \toprule
    \textbf{Requisito funcional} & \textbf{Casos de uso relacionados} & \textbf{Requisitos no funcionales relacionados} \\
    \midrule
    \endhead
    RF-01.1 & CU-01 & RNF-T-01 \\
    RF-01.2 & CU-01 & RNF-T-01, RNF-01.1 \\
    RF-01.3 & CU-01, CU-02 & RNF-01.1 \\
    RF-01.4 & CU-01 & RNF-T-05, RNF-01.1 \\
    RF-01.5 & CU-02 & RNF-T-01, RNF-01.1 \\
    RF-01.6 & CU-03 & RNF-T-01, RNF-01.1 \\
    RF-01.7 & CU-01, CU-04 & RNF-T-05 \\
    RF-01.8 & CU-02 & --- \\
    RF-02.1 & CU-04 & --- \\
    RF-02.2 & CU-05 & RNF-T-04 (compatibilidad de navegador) \\
    RF-T-01 & CU-07 & --- \\
    RF-T-02 & CU-07 & --- \\
    RF-T-03 & CU-06 & --- \\
    RF-T-04 & CU-06 & RNF-03.1 \\
    RF-T-05 & CU-08 & --- \\
    RF-T-06 & Precondición de CU-01, CU-04 & --- \\
    \midrule
    RF-03.1 & CU-09 & RNF-03.1 \\
    RF-04.1 & CU-11 & RNF-03.1 \\
    RF-06.1 & CU-10 & RNF-03.1 \\
    \bottomrule
\end{longtable}

\small
\begin{longtable}{>{\raggedright\arraybackslash}p{3.4cm} >{\raggedright\arraybackslash}p{10.2cm}}
    \caption{Matriz de trazabilidad RNF -- sub-requisitos}
    \label{tab:trazabilidad-rnf} \\
    \toprule
    \textbf{RNF} & \textbf{Sub-requisitos técnicos} \\
    \midrule
    \endfirsthead
    \toprule
    \textbf{RNF} & \textbf{Sub-requisitos técnicos} \\
    \midrule
    \endhead
    RNF-01 (latencia) & RNF-01.1 \\
    RNF-02 (precisión) & RNF-02.1 \\
    RNF-03 (seguridad/privacidad) & RNF-03.1 \\
    RNF-04 (hardware) & RNF-04.1 (compartido con RNF-07.1) \\
    RNF-05 (dataset especializado) & RNF-05.1 \\
    RNF-06 (validación lingüística) & RNF-06.1 \\
    RNF-07 (estación fija) & RNF-07.1 (compartido con RNF-04.1) \\
    \bottomrule
\end{longtable}

\subsection{Apéndice B: Formatos de datos relevantes}

\subsubsection{Formato de un frame (entrada al pipeline de reconocimiento)}
\begin{verbatim}
Frame = [
  [x_0, y_0, z_0, vis_0],     // landmark 0 (pose)
  ...                          // 33 landmarks de pose
  [x_33, y_33, z_33, vis_33],  // landmark 33 (mano izquierda)
  ...                          // 21 landmarks de mano izquierda
  [x_54, y_54, z_54, vis_54],  // landmark 54 (mano derecha)
  ...                          // 21 landmarks de mano derecha
]
// Total: 75 landmarks x 4 valores = 300 flotantes por frame
\end{verbatim}

\subsubsection{Formato de respuesta de predicción (WebSocket)}
\begin{verbatim}
{
  "sign": "dolor",
  "confidence": 88.5,
  "top3": [
    {"sign": "dolor", "confidence": 88.5},
    {"sign": "cabeza", "confidence": 8.2},
    {"sign": "idle", "confidence": 3.3}
  ]
}
\end{verbatim}

\subsubsection{Formato de respuesta de traducción (REST)}
\begin{verbatim}
{
  "original": "yo dolor cabeza",
  "natural_sentence": "Me duele la cabeza.",
  "audio_base64": "UklGRi4AAABXQVZFZm10IBAA...",
  "source": "template",
  "latency_ms": 3
}
\end{verbatim}
